\documentclass[11pt]{article}
\usepackage[margin=2.2cm]{geometry}
\usepackage{amsmath,amssymb,booktabs,enumitem,xcolor,tabularx}
\usepackage[hidelinks]{hyperref}
\setlist{nosep,leftmargin=*}
\setlength{\parindent}{0pt}
\setlength{\parskip}{4pt}

\newcommand{\nb}{\bar n}
\newcommand{\pending}[1]{\textcolor{red}{[#1]}}
\newcommand{\edit}[1]{\textcolor{blue}{[#1]}}
\newcommand{\src}[1]{{\footnotesize\texttt{#1}}}
% section header row in the claims table
\newcommand{\claimsec}[2]{\multicolumn{2}{@{}l}{\textcolor{black!60}{\S#1}\enspace\textsc{#2}} \\}

\title{Wrong Family, Right Fit: Parameter Fitting for Spatial Point Processes}
\author{AISTATS submission plan}
\date{}

\begin{document}
\maketitle
\thispagestyle{empty}

Alternative titles:
\begin{itemize}
\item Wrong Family, Right Fit: Fitting Spatial Point Processes from a Single Realization
\item Inferring Spatial Point Processes from One Realization with Topological Deep Learning
\end{itemize}

\section{Abstract -- Omer}

Spatial point processes (SPPs) are widely used to model data consisting of discrete event locations. A broad range of statistical models has been developed to capture spatial interactions, from homogeneous Poisson processes describing complete spatial independence to models capturing clustering, repulsion, and more complex spatial dependence. In many applications, only a single point pattern is available, and we seek a plausible generative model from which additional patterns with a similar underlying distribution can be simulated. This requires both selecting an appropriate model family and estimating its parameters, from a single realization. We develop a unified pipeline for these tasks across multiple SPP families, combining classical spatial summary statistics with geometric and topological features derived from \emph{persistent homology}. 
We show that these features improve estimation performance across a range of SPP families, with the resulting pipeline achieving state-of-the-art results on the evaluated benchmarks.
Beyond predictive performance, we further analyze the estimation problem itself, characterizing regimes in which parameters are difficult to recover from a single point pattern while showing that new realizations preserving the statistical properties of the observed pattern can nonetheless be generated. Together, these results provide a practical framework for fitting spatial models and generating new realizations consistent with the observed data.

%Beyond predictive performance, we identify regimes in which realizations from different distributions are statistically indistinguishable, revealing fundamental limitations of inference from a single point pattern. [[[is it a limitation? I thought the result show that for estimation it doesn't matter...]]] Together, these results provide a practical framework for fitting spatial models and generating new realizations that preserve the statistical properties of the observed pattern.

\edit{We show that these features improve estimation performance across a range of SPP families, with the resulting pipeline achieving state-of-the-art results on the evaluated benchmarks. Beyond predictive performance, we identify regimes in which realization from different distributions are statistically indistinguishable (an intrinsic property of the data, not necessarily a failure of inference) and show that this ambiguity is inconsequential in practice: even when the generating family cannot be reliably identified, our pipeline produces high-fidelity realizations of the observed pattern. Together, these results provide a practical framework for fitting spatial models and generating new realizations that preserve the statistical properties of the observed pattern.}


\section*{Abstract -- version 1}


Spatial point processes (SPPs) are  used to model data arising as collections of event locations. A wide range of models has been developed, from homogeneous Poisson processes describing complete spatial randomness to models capturing clustering, repulsion, and more complex spatial dependence. In many applications, only a single point pattern is observed, yet we seek a plausible generative model from which additional patterns with similar spatial structure can be simulated. This motivates both selecting an appropriate model family and estimating its parameters from a single realization. We develop a unified pipeline for these tasks across multiple SPP families, combining classical spatial summary statistics with geometric and topological features derived from \emph{persistent homology}. Across several model families, these features improve parameter estimation, with the resulting pipeline achieving state-of-the-art performance on the evaluated benchmarks. Beyond predictive performance, we analyze the conditions under which estimation becomes difficult, providing insight into the challenges of inference from a single point pattern. Together, these results provide a practical approach to fitting plausible spatial models and generating additional realizations that reflect the structure of the observed data.

\section{abstract -- version 2}
Spatial point processes (SPPs) are commonly used to model data arising as collections of event locations. A wide range of models has been developed, from homogeneous Poisson processes describing complete spatial randomness to models capturing clustering, repulsion, and more complex spatial dependence. In many applications, only a single point pattern is observed, but we require a plausible generative model from which similar spatial patterns can be simulated from an established parametric SPP model. In this paper, we study model classification and parameter estimation in this setting, developing a pipeline that combines classical spatial summary statistics with geometric and topological features derived from \emph{persistent homology}. Across several SPP families, these features improve parameter estimation, with the resulting pipeline achieving state-of-the-art performance. Finally, we also provide an in-depth analysis of the estimation problem, identifying a class of parameters that are difficult to estimate and showing how these difficulties can be mitigated.

\section{Abstract -- short}
Spatial point processes (SPPs) are common models for random event locations. Given a single observed pattern, we aim to find a plausible model from an established parametric SPP family to generate similar patterns. We develop a pipeline combining classical spatial summary statistics with topological features derived from \emph{persistent homology} for model classification and parameter estimation. These features improve parameter estimation across several SPP families, yielding state-of-the-art performance. 
We further analyze the estimation problem itself, characterizing a regime in which parameters are difficult to recover from a single pattern and showing how plausible point patterns can be nonetheless produced. 

%Beyond the pipeline, we perform an in-depth analysis of the estimation problem, identifying a class of problems in which it is difficult to recover parameters from a single observed pattern and showing how estimation can be improved.

\section{Abstract -- old}
\vspace{2cm}

Spatial point processes (SPPs) are widely used to model real-world phenomena, and their parameters can be directly related to physical quantities of interest.

The quality of an SPP fit is typically assessed by the distance between inferred and true generating parameters. Existing methods, such as minimum contrast, require prior knowledge of the underlying family, and their performance degrades near complete spatial randomness (CSR): despite the generating process not being Poisson, its second-order structure is almost indistinguishable from Poisson's. This is an inherent property of the problem, not a failure of methodology.

Often, practitioners simply need high-fidelity replicates of the original clouds. We therefore score a fit by how well realizations generated with the predicted parameters reproduce an independent held-out replicate, using a strictly proper scoring rule on local configurations.

We propose a two-stage pipeline: a classifier determines the family, followed by family-specific parameter estimators, with the first stage acting as a filter that routes near-CSR patterns to the Poisson estimator. Both stages draw on second-order summary functions and topological features from persistent homology. 

Near CSR, patterns become inherently difficult to classify regardless of the model used. We show that in these cases naming the correct family matters little: the pipeline generates high-fidelity replicates even though the patterns are often misclassified. Topological features never hurt parameter estimation, and help substantially where second-order statistics are blind by construction (e.g. cell process).

%==============================================================================

\section*{The Story [INFORMAL]}

\begin{itemize}
    \item All models struggle with classification/inference on clouds near CSR. Unless you cherry-pick your prior (bad), you will have to deal with a large number of near-CSR clouds. Bad performance on these near-CSR clouds is not a failure of the models, it's an inherent property of the problem. 

    \item The pitch: we show that, despite this difficulty, classifiers successfully tell us which generating process is "good enough" for inference.

    \item Classification accuracy on the 8 processes seems weak ($0.53$), but the final fidelity score - i.e. how well the inferred parameters reproduce the pattern - is high ($0.92$).

    \item We show that when near-CSR but non-Poisson clouds get routed to Poisson inference, the resulting fit is still high-fidelity.

    \item The classical alternative (minimum contrast on $K$) does worse at every stage: selection, estimation, and end-to-end fidelity.
\end{itemize}

\newpage

%==============================================================================

\section*{Contributions}

\begin{enumerate}
  \item \textbf{Scoring.} We evaluate a fit by whether patterns simulated from it look like a held-out replicate with the same parameters (for each $\theta$ in the cloud bank we generate two realizations, one of which is held-out for scoring).
  \begin{itemize}
    \item \emph{General idea.} $P_j = (\hat f, \hat \theta)$ is the fit for cloud $x_j$ and is scored on the independent replicate $y_j$ (produced with the same $\theta$). We also compute scores for two reference models on $y_j$: the oracle $P^\star_j$ (true family and $\theta$) and Poisson $P^0_j$ with $\bar n=|y_j|$.
    \item \emph{Local configurations.} Rescale a pattern to unit intensity (multiplying points by $\sqrt{n}$). Around a number of points (up to $m=300$), take neighbours within $R=1.5$ mean spacings, translated to the origin and divided by $R$.
    \item \emph{Kernel.} Embed each local configuration $\phi$ as a smoothed density on a $21\times21$
      grid:
      \[
      \mu_\phi(u)=\sum_{p\in\phi}e^{-|u-p|^2/2h^2}, \qquad h = 0.1
      \]
    Compare configurations with 
    \[
    k(\phi,\psi)=\exp\!\big(-\|\mu_\phi-\mu_\psi\|^2/2\ell^2\big)
    \]
    where $\ell$ is the median distance between $y$'s own configurations, fixed before any model is scored.
    \item \emph{Score} (lower is better):
      \[
        S(P,y)=\tfrac12\,\mathbb{E}\,k(\Phi,\Phi')-\tfrac1m\textstyle\sum_{i=1}^m\mathbb{E}\,k(\phi_i,\Phi),
      \]
      where $\phi_i$ are $y$'s configurations and $\Phi,\Phi'$ are configurations from different
      simulations of $P$.
    \item \emph{Skill}, pooled over test patterns and reported only where the denominator is
      resolved ($>2$ s.e.):
      \[
        \text{skill}=\frac{\sum_j\big[S(P_{0,j},y_j)-S(\hat P_j,y_j)\big]}{\sum_j\big[S(P_{0,j},y_j)-S(P^\star_j,y_j)\big]},
      \]
      so $0$ means no better than CSR and $1$ means as good as the truth.
  \end{itemize}
  \item \textbf{CSR threshold.} For every family except cell (whose $K$ is exactly Poisson's; there the number of points per cell $k$ plays the role) we show that whether a pattern can be distinguished from CSR is governed by a physical coordinate (cluster overlap $\omega$, or a pair-count signal-to-noise ratio $\zeta$) multiplied by a fitted power of $\nb$. Past this boundary all classifiers and estimators improve; before it the proper score shows that even the oracle (i.e. true family with true parameters) does not beat Poisson inference.
  \item \textbf{Pipeline.} Classifier followed by family-specific inference model, reflecting the factorization $p(f,\theta\mid x)=p(f\mid x)\,p(\theta\mid x,f)$. The classifier's Poisson call acts as the CSR filter: structured patterns it sends to Poisson cost nothing, and patterns it assigns to the wrong structured family are reproduced about as well as correctly identified ones. The pipeline beats the classical minimum-contrast fit at selection, estimation and end-to-end fidelity.
  \item \textbf{TDA features.} Adding PH to classical summaries never makes estimates significantly worse and helps significantly where second-order statistics cannot see the structure (e.g. cell process which has $K(r)=\pi r^2$ exactly).
\end{enumerate}

\newpage

%==============================================================================

\section*{Paper outline}

\begin{tabularx}{\linewidth}{@{}p{2.8cm}X@{}}
\toprule
Section & Content \\
\midrule

1 Introduction &
Hook: accuracy $0.53$ vs skill $0.92$. Why accuracy is the wrong target near CSR. Contributions. \\

2 Setting \& score &
Introduce families and data generation process (50\,000 $\theta$ $\times$ 2 replicates per family). Kernel score, anchors, skill.\\

3 Model &
Two-step pipeline: classifier as filter, then per-family estimators on $\log\theta$. Inputs: 111 classical summary-function samples ($L,F,G,J$), 105 PH summaries per filtration, PersLay, persistence images. Learners: gradient-boosted trees, linear, CNNs. \\

4 CSR wall &
Regime coordinates directly explain difficulty: show per-model detection curves, parameter error vs coordinate, and oracle-vs-CSR gain vs coordinate. For non-Poisson clouds with CSR-like structure, calling them Poisson costs nothing. \\

5 Fidelity &
Report overall results and analysis. Show that classifier performance shoots up when CSR clouds are omitted. Provide specific metrics on misclassified clouds, proving that whether they are inferred with the right family model or not doesn't change fidelity. Compare end to end against minimum contrast (selection and estimation). \\

6 TDA &
Feature ablation table, show that TDA either does nothing (i.e. doesn't actively hurt), or significantly boosts performance: might as well include always. \\

\bottomrule
\end{tabularx}

\newpage

%==============================================================================

\section*{Main numbers behind each claim}

{\small
\begin{itemize}
    \item Test split. Accuracy is balanced over the 8 families.
    \item Estimation error: RMSE$(\log\theta)$/s.d., lower is better.
    \item Skill: 1 = as good as the true model, 0 = no better than CSR.
    \item In regime at $\tau$: structured patterns with $P(\text{detected})\ge\tau$.
    \item Pipeline: one architecture (HGB on classical + PH) for the classifier and every per-family estimator;
    \item Classical baseline: min contrast on $K$.
\end{itemize}

\medskip 

{\renewcommand{\arraystretch}{1.45}
\begin{tabularx}{\linewidth}{@{}p{7.2cm}X@{}}
\toprule
\textbf{Claim} & \textbf{Evidence} \\
\midrule
\claimsec{1}{Introduction}
Accuracy looks weak, fidelity is high & Accuracy $0.53$ vs. skill $0.92$ \\
Classification errors cost little fidelity & Skill $0.92$ with predicted family vs $0.94$ with correct one \\
\midrule
\claimsec{4}{The CSR wall}
Accuracy has a ceiling & Best classifier $0.535$, all 11 combined $0.537$ \\
One coordinate governs detection & Detection AUC $0.87$--$0.96$, vs $0.88$--$0.96$ from all parameters \\
Models improve past the boundary & Accuracy $0.53\to0.65$, Thomas error $-37\%$ \\
Before the boundary, even the true family does not beat CSR & Oracle beats CSR in 0 of 14 strata below the boundary \\
Routing to Poisson is free & Oracle gain on structured patterns sent to Poisson $\approx0$ \\
\midrule
\claimsec{5}{Fidelity}
Pipeline reproduces patterns & Skill $0.92$ vs $0.84$ for min.\ contrast \\
Wrong family, right fit & Skill $0.91$ misclassified vs $0.95$ identified (difference n.s.) \\
Pipeline beats min.\ contrast & Selection $0.53$ vs $0.33$; min.\ contrast Thomas error $2.3\times$ higher \\
Min.\ contrast degrades near CSR & Thomas error $3.3\times$ higher outside the regime \\
\midrule
\claimsec{6}{TDA}
PH never hurts & No estimator significantly worse with PH \\
PH helps where $K$ is blind & Cell error $-15\%$ \\
\ldots\ and it shows end to end & Cell skill $0.92\to1.00$ \\
\bottomrule
\end{tabularx}}

\newpage

%==============================================================================

\section*{Figures and tables}

{

\medskip
{\renewcommand{\arraystretch}{1.45}
\begin{tabularx}{\linewidth}{@{}lp{1.2cm}Xl@{}}
\toprule
ID & Section & Content & Status \\
\midrule
\multicolumn{4}{@{}l}{\textbf{MAIN TEXT}} \\
T1 & 2 & The 8 families: parameters, $K$, regime coordinate & written \\
F1 & 2 & Example realizations, one strongly structured and one near CSR per family & generated \\
F2 & 3 & Pipeline schematic: cloud $\to$ classifier $\to$ estimator $\to$ score & \pending{to do} \\
F4 & 4 & $P$(called Poisson) over each family's parameter space, with the fitted boundaries & generated \\
F5 & 4 & \emph{Key figure}: detection, estimation error and oracle gain over CSR against the regime coordinate, one curve per model & generated \\
T6 & 4 & Oracle and fitted-model gain over CSR by regime stratum; ``structured, called Poisson'' row & generated \\
T7 & 5 & Main results: accuracy, error, skill for pipeline / classical / min.\ contrast / true-family routing; wrong-family rows & \pending{to build} \\
T10 & 6 & What PH adds: paired estimator differences, classifier accuracy, skill with and without & generated \\
\multicolumn{4}{@{}l}{\textbf{APPENDIX}} \\
T2 & 2 & Simulation design and data volume & written \\
T3 & 3 & Models: learner, input, training cost & generated \\
T8 & 2 & Score power: $d'$ of true model vs CSR, kernel vs alternatives & \pending{rerun} \\
T4 & 4 & Classification by model $\times$ family, overall and in regime, with min.\ contrast & generated \\
F3 & 4 & Confusion matrix of the frozen classifier & generated \\
T5 & 4 & Fitted ``looks like Poisson'' boundary per family & generated \\
T9 & 5 & Estimation error per family $\times$ estimator, overall and in regime, with min.\ contrast & generated \\
\bottomrule
\end{tabularx}}
}

\end{document}
